\par\Needspace{4\baselineskip}
\begin{center}
\SourceLeaf{022b}\pdfbookmark[2]{4.}{article-004}4.
\end{center}
In place of the distance $x$, let us now introduce the angle $v$ between the line of apsides and the straight line drawn from the Sun to the position of the body (the \emph{radius vector}). Let this angle be measured from the part of the line of apsides where the distances $x$ are positive, and let it increase in the direction of the body's motion. Thus $x=r\cos v$, and our formula becomes $r=\frac{p}{1+e\cos v}$. We immediately obtain the following conclusions.

I. For $v=0$ the radius vector $r$ attains its minimum value, namely $=\frac{p}{1+e}$. This point is called the \emph{perihelion}.

II. Opposite values of $v$ correspond to equal values of $r$. The line of apsides therefore divides the conic section into two equal parts.

III. In the \emph{ellipse}, $r$ increases continuously from $v=0$ until it reaches its maximum value $\frac{p}{1-e}$ at the \emph{aphelion} for $v=180^{0}$. Beyond aphelion it decreases again in just the way it had increased before, until perihelion is reached again at $v=360^{0}$. The part of the line of apsides bounded on one side by perihelion and on the other by aphelion is called the \emph{major axis}. Thus the semimajor axis, also called the \emph{mean distance}, is $=\frac{p}{1-ee}$. The distance of the midpoint of the axis (the \emph{centre of the ellipse}) from the focus is $\frac{ep}{1-ee}=ea$, with $a$ denoting the semimajor axis.

IV. In the \emph{parabola}, on the other hand, there is strictly speaking no aphelion. The value of $r$ increases without bound as $v$ approaches $+180^{0}$ or $-180^{0}$. For $v=\pm180^{0}$ the value of $r$ becomes infinite, which shows that the line of apsides does not intersect the curve on the side opposite perihelion. Thus, strictly speaking, the curve has neither a major axis nor a centre. In accordance with the usual practice of analysis, however, we extend the formulas found for the ellipse and assign an infinite value to the major axis, placing the centre of the curve at an infinite distance from the focus.

\SourceLeaf{023}V. Finally, in the \emph{hyperbola}, $v$ is confined within still narrower limits, namely between $v = -(180^0-\psi)$ and $v = +(180^0-\psi)$, where $\psi$ denotes the angle with cosine $= \frac{1}{e}\cdot$ Indeed, as $v$ approaches these limits, $r$ increases without bound. If either limit itself were taken for $v$, the value of $r$ would be infinite, which shows that a straight line inclined at an angle of $180^0-\psi$ above or below the line of apsides does not intersect the hyperbola at all. For the values thus excluded, namely from $180^0-\psi$ to $180^0+\psi$, our formula assigns a negative value to $r$. A straight line inclined at such an angle to the line of apsides does not itself intersect the hyperbola, but if extended backwards it meets the other branch. This branch is entirely separate from the first and is known to be convex towards the focus occupied by the Sun. In our investigation, however, which rests on the assumption already stated that $r$ is taken as positive, we shall disregard this other branch of the hyperbola. Only a celestial body on which the Sun exerted a repulsive force according to the same laws, instead of an attractive force, could travel along it. — Strictly speaking, then, there is no aphelion in the hyperbola either. The point of the opposite branch that lies on the line of apsides and corresponds to $v = 180^0$, $r = -\frac{p}{e-1}$ may be regarded as analogous to the aphelion. If, by analogy with the ellipse, we choose to call the value of the expression $\frac{p}{1-ee}$ the semimajor axis of the hyperbola even here, where it is negative, then twice this quantity gives the distance of the point just mentioned from perihelion and indicates that it lies on the opposite side of perihelion from the corresponding point in the ellipse. Similarly, $\frac{ep}{1-ee}$, that is, the distance from the focus to the midpoint between these two points (the centre of the hyperbola), takes a negative value here because of its opposite position.
